% !TeX root = ../main.tex

\chapter{补充材料}

\section{实验复现文件说明}

本文在定稿阶段保留了完整的实验目录，核心文件如下。

\begin{itemize}
  \item \path{code/data/skills_full.json}：完整数据资产，包含 10028 条 Skill 记录。
  \item \path{code/results/nightly-20260330-v2/}：1000 规模主实验结果目录，内含 \path{summary.json}、\path{artifacts/1000/terrain.json} 与各类日志文件。
  \item \path{code/results/scale-100-tuned/}：100 规模独立调优结果目录。
  \item \path{code/results/thesis-ready-summary-20260330.json}：用于论文回填的整理版摘要文件。
  \item \path{image/chap04/demo-home-3d.png} 与 \path{image/chap04/demo-query-3d-test.png}：论文中直接使用的 3D Demo 截图。
\end{itemize}

\section{复现实验命令}

以下命令可在当前仓库中直接复现实验与 Demo：

\begin{lstlisting}[language=bash]
code/.venv/bin/python code/scripts/run_experiments.py \
  --sizes 100,1000 \
  --output-dir code/results/nightly-20260330-v2

code/.venv/bin/python code/scripts/serve_demo.py \
  --results-dir code/results/nightly-20260330-v2 \
  --dataset-size 1000 \
  --port 8765
\end{lstlisting}

\section{日志与后续迭代建议}

\begin{itemize}
  \item 终端原始输出分别保存在 \path{code/results/nightly-20260330-v2-run.log} 与 \path{code/results/scale-100-tuned-run.log} 中，可直接用于回看每轮执行细节。
  \item AutoResearch 的调参轨迹保存在 \path{logs/autoresearch_trace.json} 中，适合后续继续沿同一参数搜索空间做增量实验。
  \item 若下一轮需要补跑更大规模实验，建议在当前目录结构下新建独立结果目录，避免覆盖本轮论文定稿所依赖的产物。
\end{itemize}
